\chapter{Introduction}\label{chapter:introduction}
\thispagestyle{myheadings}

	% ========================================================================
	% GENERAL INTRODUCTION
	%
	% In a life-science dissertation this chapter does the work that a review
	% article would: it establishes the field, narrows to the open question,
	% and states what you did about it. Aim for 15-30 pages.
	%
	% A conventional structure, which the sections below follow:
	%   1. Broad significance -- why should anyone care?
	%   2. Current state of knowledge -- what is established?
	%   3. The gap -- what is missing or contradictory?
	%   4. Rationale for the approach -- why this technique, this model?
	%   5. Specific aims -- one paragraph per results chapter.
	% ========================================================================

	TODO: Open with a paragraph that a scientist outside your subfield could
	follow, stating the broad problem and its clinical or biological stakes.

	\section{Background}\label{section:introduction:background}

		TODO: Review the established literature. Build the argument in the
		order the reader needs it, not in the order you learned it.

		\subsection{TODO: first background theme}

			TODO: Cite with \verb|\parencite{key}| for a parenthetical
			citation, \verb|\textcite{key}| when the authors are the subject
			of your sentence, and \verb|\cite{key}| for the plain form. For
			example: hemodynamic coupling has been characterised extensively
			\parencite{placeholder-review}, and \textcite{placeholder-methods}
			introduced the measurement approach used here.

		\subsection{TODO: second background theme}

			TODO.

	\section{Gap in current knowledge}\label{section:introduction:gap}

		TODO: State plainly what is not known. This paragraph is the hinge of
		the whole dissertation -- everything before it motivates the gap, and
		everything after it addresses the gap.

	\section{Rationale and approach}\label{section:introduction:rationale}

		TODO: Justify your model system, your measurement modality and your
		manipulation. Explain what each one can and cannot resolve. Anticipate
		the obvious "why didn't you just do X instead" question here rather
		than in the discussion.

	\section{Specific aims}\label{section:introduction:aims}

		\begin{description}[style=unboxed, leftmargin=0em]

			\item[Aim 1.]
				TODO: One or two sentences. Addressed in
				\cref{chapter:study-one}.

			\item[Aim 2.]
				TODO: One or two sentences. Addressed in
				\cref{chapter:study-two}.

		\end{description}

	\section{Organization of this dissertation}\label{section:introduction:structure}

		\Cref{chapter:methods} describes the animals, surgical preparations,
		optical instrumentation, experimental designs and analysis pipelines
		shared by all of the studies that follow.
		\Cref{chapter:study-one} reports TODO.
		\Cref{chapter:study-two} reports TODO.
		\Cref{chapter:discussion} integrates the findings, states the
		limitations and outlines future directions.

	\subsection*{Work completed during the doctoral program}

		% If a chapter has been published, you MUST say so on the copyright
		% page as well (BU guide 1.3.2), e.g.:
		%   (c)2027 by Daria Bogatova. All rights reserved except for
		%   chapter 3, which is (c)2026 Journal Name.
		% List your publications here with \printpublication{citekey}.
		TODO: List the papers, preprints and conference abstracts that came
		out of this work, and note which chapters they correspond to.
